\documentclass[11pt,a4paper]{article}
\usepackage[margin=0.85in,top=1in,bottom=1in]{geometry}
\usepackage{amsmath,amssymb,amsfonts}
\usepackage{graphicx}
\usepackage{float}
\usepackage{enumitem}
\usepackage{microtype}
\usepackage{caption}
\usepackage[hidelinks]{hyperref}
\usepackage[most,breakable]{tcolorbox}
\tcbuselibrary{breakable,skins}
\usepackage[utf8]{inputenc}
\usepackage[T1]{fontenc}
\usepackage{lmodern}
\captionsetup{font=small,labelfont=bf,skip=4pt}
\setlist{leftmargin=1.5em,topsep=2pt}

%% Breakable card — prevents content cutoff across pages
\newtcolorbox{solcard}[3]{
  breakable, enhanced,
  colback=white, colframe=gray!50, arc=2pt,
  title={\textbf{#1}\hfill\textsf{\small #3\,pts}},
  fonttitle=\normalsize\sffamily\bfseries,
  before upper={\small\textit{#2}\par\smallskip\hrule height 0.4pt\smallskip},
  top=4pt, bottom=6pt, left=7pt, right=7pt,
  parbox=false,
}

\title{\textbf{In the middle of the bridge --- Solution}}
\author{\normalsize W24 (2024) — English translation}
\date{}
\begin{document}\maketitle\vspace{-0.5em}
\begin{solcard}{A1}{If the section of the circuit between nodes $A$ and $B$ is removed, the remaining electrical circuit consists entirely of constant voltage sources and resistors. Therefore, the dependence of the current $I_{AB}$ flowing from node $A$ to node $B$ on the voltage $U_{AB}=\varphi_A-\varphi_B$ across this section can be represented in a form that does not depend on the specific element connected between nodes $A$ and $B$: $$I_{AB}=\cfrac{\mathcal{E}_0-U_{AB}}{r_0}{.} $$Determine $\mathcal{E}_0$ and $r_0$. Express your answer in terms of $\mathcal{E}$ and $R$.}{1.50}

If an ideal voltmeter is connected between nodes $A$ and $B$, the voltage across it is $$ U_V = \mathcal{E}\left( \frac{6R}{6R + 4 R} - \frac{2 R}{2 R + 8 R}\right) = \frac{2}{5} \mathcal{E} = \mathcal{E}_0. $$ If the voltage of the source is set to zero, it can be replaced by a wire with zero resistance. Then the circuit diagram between terminals $A$ and $B$ can be represented as two pairs of parallel-connected resistors, $(2R,\, 8R)$ and $(6R,\, 4R)$, connected in series, which means the equivalent resistance is $$ r_0 = R_{AB} = \frac{2R \cdot 8R}{2 R + 8 R} + \frac{6R \cdot 4R}{6R + 4 R} = 4R. $$Then the dependence of the voltage on the current has the form $$ U_{AB} = \mathcal{E}_0 - I_{AB} r_0, $$from which $$ I_{AB} = \frac{\mathcal{E}_0 - U_{AB}}{r_0}. $$Alternatively, one can calculate the short-circuit current $I_c = \mathcal{E}_0/r_0$, which would flow through an ideal ammeter connected between terminals $A$ and $B$. For such a connection, the total resistance of the circuit to which the source $\mathcal{E}$ is connected equals $$ r_1 = \frac{2 R \cdot 6 R}{2 R + 6 R} + \frac{8 R \cdot 4 R}{8R + 4 R} = \frac{25}{6} R. $$Then the total current through the source $\mathcal{E}$ is $I_\mathcal{E} = 6\mathcal{E}/25 R$, and the currents through the $2R$ and $8R$ resistors are respectively: $$ I_1 = \frac{3}{4} I_\mathcal{E} = \frac{9\mathcal{E}}{50 R}, \quad I_2 = \frac{1}{3}  I_\mathcal{E}  = \frac{2\mathcal{E}}{25 R}. $$Then the short-circuit current is $$ I_c = I_{AB} = I_1 - I_2 = \frac{\mathcal{E}}{10 R}. $$Hence, we obtain the previous value for the resistance: $$ r_0 = \mathcal{E}_0/I_c = 4 R. $$

\par\smallskip\noindent\textbf{Answer:} $$r_0 = 4 R, \quad \mathcal{E}_0 = \frac{2}{5} \mathcal{E} $$
\end{solcard}

\begin{solcard}{A2}{Obtain the dependence of the current strength $I_{AB}$ on the value $\mathcal{E}$. Express your answer using $\mathcal{E}$, $R$, $U_0$ and $I_0$.}{0.50}

Suppose that the voltage across the nonlinear element is less than $U_0$. Then this voltage can be found from the equation $$ (r_0 + R) I_{AB} = \mathcal{E}_0 - U_{AB}, \quad (r_0 + R) I_0 \frac{U_{AB}}{U_0} = \mathcal{E}_0 - U_{AB}, $$from which $$ U_{AB} =  \frac{\mathcal{E}_0 U_0}{U_0 + (r_0 + R) I_0}, \quad \mathcal{E}_0 \le U_0 + (r_0 +R) I_0. $$The corresponding value of the current is $$ I_{AB} =  \frac{\mathcal{E}_0 I_0}{U_0 + 5 R I_0}, \quad \mathcal{E}_0 \le U_0 + 5R I_0. $$With a further increase in the source voltage, the voltage across the nonlinear element becomes greater than $U_0$, and the current through it equals $I_0$. The voltage across the nonlinear element is $U_{AB} = \mathcal{E}_0 - I_0 r_0$. Substituting the expressions for $\mathcal{E}_0$ and $r_0$, we get the answer.

\par\smallskip\noindent\textbf{Answer:} $$ I_{AB} = \begin{cases} \frac{2 \mathcal{E} I_0}{5 (U_0 + 5R I_0)},  \quad \mathcal{E} \le \frac{5}{2} \left(U_0  + 5 R I_0 \right);\\ I_0, \quad \mathcal{E} > \frac{5}{2} \left(U_0  + 5 R I_0 \right). \end{cases} $$
\end{solcard}

\begin{solcard}{A3}{Obtain the dependence of the power $P_\text{ne}$ dissipated by the nonlinear element on the value $\mathcal{E}$. Express your answer using $\mathcal{E}$, $R$, $U_0$ and $I_0$.}{0.50}

The power dissipated by the nonlinear element is $P_{\text{ne}} = U_{\text{ne}} I_{AB}$. Using the results from the previous section, we get the answer.

\par\smallskip\noindent\textbf{Answer:} $$ P_{\text{ne}} = \begin{cases} U_0 I_0 \cfrac{4 \mathcal{E}^2}{25 (U_0 + 5 R I_0)^2}, \quad  \mathcal{E} \le \frac{5}{2} \left(U_0  + 5 R I_0 \right);\\ I_0 \left( \frac{2}{5} \mathcal{E} - 5 I_0 R\right),  \quad  \mathcal{E} > \frac{5}{2} \left(U_0  + 5R I_0 \right). \end{cases} $$
\end{solcard}

\begin{solcard}{A4}{Construct a high-quality graph of the dependence $P_\text{ne}(\mathcal{E})$. Indicate all the characteristic points on the graph and indicate their corresponding values.}{0.50}

At low voltages, the graph is a parabola with its vertex at the origin. At $\mathcal{E} > \frac{5}{2} \left(U_0  + 4 R I_0 \right)$ it becomes a straight line. At the point where the transition from a quadratic to a linear dependence occurs, $\mathcal{E} = U_0 + 4 R I_0$ and $P_{\text{ne}} =  I_0 U_0$.

\begin{figure}[H]
  \centering
  \includegraphics[width=0.72\textwidth]{T2_S_EN_assets/fig1.jpg}
\end{figure}
\end{solcard}

\begin{solcard}{B1}{Obtain the dependence of the current $I_C(t)$ through the capacitor on time $t$. Express your answer in terms of $\mathcal{E}$, $R$, $U_0$, $I_0$, $C$ and $t$.}{0.80}

If the voltage across a nonlinear element is less than $U_0$, it behaves like a resistor with a resistance of $r_0 = U_0/I_0 = 500~\mathrm{\Omega} = R$. Therefore, in this case, the voltage across the nonlinear element must be equal to the voltage across the resistor. Then the initial voltage on the nonlinear element is equal to $\mathcal{E}/2  = 3~\mathrm{V}< U_0$, so we can use the linear dependence of current on voltage everywhere. Then the charge of the capacitor and the current through it are related by the following relation (which is obtained from Kirchhoff's equation): $$ I (R + r_0) + \frac{q}{C} = \mathcal{E}. $$ Differentiating it with respect to time and noting that $\dot{q} = I$, we obtain $$ \dot{I} = -\frac{I}{C (R + r_0)}, \quad I(t) = I(0) e^{- \frac{t}{C(R + r_0)}}. $$At the initial moment in time, the capacitor is uncharged; therefore, $I(0) = \mathcal{E}/(R + r_0)$. Then the current through the capacitor is:

\par\smallskip\noindent\textbf{Answer:} $$ I_C = \frac{\mathcal{E}}{R + U_0/I_0} e^{- \frac{t}{C(R + U_0/I_0)}} $$
\end{solcard}

\begin{solcard}{B2}{What amount of heat $Q_\text{ne}$ will be released on the nonlinear element during the establishment of a constant current? Express your answer in terms of $\mathcal{E}$, $R$, $U_0$, $I_0$ and $C$. Find the numerical value.}{0.60}

When charging the capacitor, a charge $q = C\mathcal{E}$ passes through the source, and it does work $A_{\text{source}}  = q \mathcal{E} = C \mathcal{E}^2$. This work goes into increasing the energy of the capacitor and the heat released: $$ A_{\text{source}}  = \frac{C \mathcal{E}^2}{2} + Q. $$Therefore, the total heat released in the circuit is $Q = C \mathcal{E}^2/2$. It is divided between the resistor and the nonlinear element in proportion to their resistances (since, at the voltages in question, the nonlinear element behaves like a resistor):

$$ Q_{\text{ne}} = \frac{r_0}{r_0 + R} \frac{C \mathcal{E}^2}{2} = \frac{U_0}{U_0 + R I_0} \frac{C \mathcal{E}^2}{2} = 0.9~\mathrm{mJ} $$

\end{solcard}

\begin{solcard}{B3}{After what time $\tau$ will the amount of heat released on the nonlinear element reach $80\%$ of the value $Q_\text{ne}$? Express your answer in terms of $R$, $U_0$, $I_0$, and $C$. Find the numerical value.}{0.60}

The power on the nonlinear element as a function of time has the form: $$ P = I_C^2 \frac{U_0}{I_0} = A e^{- 2t/t_0}, \quad A = \frac{\mathcal{E}^2 U_0 I_0}{(R I_0 + U_0)^2}, \quad t_0 = C(R + U_0/I_0). $$Then over time $t$, the heat released on the nonlinear element is $$ Q(t) = \int_0 ^t P(t) dt = A \frac{t_0}{2} \left( 1 - e^{- 2t/t_0}\right) = Q_{\text{ne}} \left( 1 - e^{- 2t/t_0}\right). $$At the moment in time when $Q = 0.8 Q_{\text{ne}}$, we have $e^{-2t/t_0} = 0.2$.

\par\smallskip\noindent\textbf{Answer:} $$ t = \frac{\ln 5}{2} t_0 =  \frac{\ln 5}{2} C \left(R + \frac{U_0}{I_0}\right) = 0.0804~\mathrm{s} $$
\end{solcard}

\begin{solcard}{B4}{Obtain the dependence of the current $I_C(t)$ through the capacitor on time $t$. Express your answer in terms of $\mathcal{E}$, $R$, $U_0$, $I_0$, $C$ and $t$.}{1.50}

Now the initial voltage on the nonlinear element is greater than $U_0$, which means the current through the nonlinear element (which is equal to the current through the capacitor) equals $I_C = I_0$. This current will remain constant until the voltage across the nonlinear element drops to $U_0$. Up to this point, the charge of the capacitor depends on time as $q = I_0 t$, which means the voltage across the nonlinear element is $$ U_{\text{ne}}  = \mathcal{E} - R I_0 -I_0 t/C. $$This dependence is valid while $U_{\text{ne}} \ge U_0$, that is: $$ t \le t_1 = \frac{C}{I_0} (\mathcal{E}   - U_0 - R I_0 ). $$After this, the current will decay exponentially, just as in an ordinary charging capacitor: $$ I = I_0 \exp\left(- \dfrac{t- t_1}{C(R + U_0/I_0)}\right). $$

\par\smallskip\noindent\textbf{Answer:} $$ I_C = \begin{cases} I_0, \quad t \le \frac{C}{I_0} (\mathcal{E} - U_0 - R I_0 );\\ I_0 \exp\left(- \dfrac{t- C(\mathcal{E}-U_0 - R I_0)/I_0}{C(R + U_0/I_0)}\right), \quad t >\frac{C}{I_0} (\mathcal{E}  -U_0- R I_0 ). \end{cases} $$
\end{solcard}

\begin{solcard}{B5}{What amount of heat $Q_\text{ne}$ will be released on the nonlinear element during the establishment of a constant current? Find the numerical value with an accuracy of at least $1\%$.}{1.50}

While the current remains constant, the power dissipated by the nonlinear element is $$ P = I_0 U_{\text{ne}} = I_0 (\mathcal{E}  - R I_0 -I_0 t/C). $$Then over the time $t_1$, the released heat will be: $$ Q_1 = \int_0^{t_1} dt I_0 (\mathcal{E}  - R I_0 -I_0 t/C) = I_0 (\mathcal{E} - R I_0) t_1 - I_0 ^2t_1^2/2C, $$$$ Q_1 = C(\mathcal{E} - R I_0)(\mathcal{E} -U_0- R I_0) -  C(\mathcal{E}-U_0 - R I_0)^2/2 = \frac{C}{2} \left( (\mathcal{E} - I_0 R)^2  - U_0^2\right) = 1.2~\mathrm{mJ} $$The further dependence of power on time has the form ($t' = t - t_1$): $$ P = \frac{U_0}{I_0} I_0^2 e^{- 2t'/t_0}, \quad t_0 = C (R + U_0/I_0) $$which means the total released heat from this phase is: $$ Q_2 = \int_0 ^\infty dt' U_0 I_0  e^{- 2t'/t_0} = \frac{1}{2} U_0 I_0 t_0 = \frac{1}{2} U_0 I_0 C (R + U_0/I_0) = 2.5~\mathrm{mJ} $$

\par\smallskip\noindent\textbf{Answer:} $$Q = Q_1 + Q_2  = \frac{C}{2} \left( (\mathcal{E} - I_0 R)^2  - U_0^2\right) + \frac{C}{2} U_0 (U_0 + I_0 R)  =  \frac{C}{2} \left((\mathcal{E} - I_0 R)^2 + U_0 I_0 R \right) = 3.7~\mathrm{mJ}$$
\end{solcard}

\end{document}